\scanpage{56}
injectif\unsure{)|]} \deleted{\unsure{-- ou encore que}}.

\uline{Prop} Ce foncteur
\diagram{\begin{tikzcd}[ampersand replacement=\&,column sep=large]
(\text{Contours \unsure{comb.|comb}})\arrow[r]\& G_1\text{-ens }\text{``admissibles''}
\end{tikzcd}}
est une équivalence des catégories.
C'est une tautologie\ldots.

\margin{\uline{NB} la condition $\sigma_0x\ne\sigma_1x\quad\forall x\in R$ signifie que l'on exclut les ``monogones'' = polygones à 1 côté\ldots}
\textcircled{\small 3} \deleted{Pour le voir, \unsure{on}} \inserted{Il est utile d'}écrire $G_1$ \inserted{$=D_\infty$} à l'aide des
générateurs \deleted{$\unsure{\sigma_0,\sigma_1|\sigma,\sigma_1|\sigma_1}$ et}
\[
\left\{\begin{aligned}\sigma&=\sigma_1\\u&=\deleted{\unsure{\sigma_2|\sigma_1}}\sigma_1\sigma_0\end{aligned}\right.
\quad\Longrightarrow\quad
\begin{aligned}\sigma_1&=\sigma\\\sigma_0&=\deleted{\unsure{u|1|\sigma}}\sigma u=u^{-1}\sigma\end{aligned}
\]
alors on a
\[
G_1=\{\sigma,u\mid\sigma^2=1,\ \deleted{\unsure{u}}\sigma u\sigma^{-1}=u^{-1}\}
\]
i.e.
\[
G_1=\underbrace{\mathbb Z.\mathbb Z/2\mathbb Z}_{\substack{\Vert\ \mathrm{dfn}\\D_\infty\\\text{(groupe diédral infini)}}}
\]
\margin{produit $\frac12$ direct, avec $\sigma$ (gén. de $\mathbb Z/2\mathbb Z$) opérant sur $\mathbb Z$ par symétrie}

La condition que \inserted{les involutions} $\sigma_0$ et $\sigma_1$ \deleted{\unsure{op}} opèrent ss pt
fixe dans $R$ s'écrivent par
\[
\left\{\begin{aligned}\sigma x&\ne x\\u^{-1}x&\ne\sigma x\end{aligned}\right.\qquad\forall x\in R
\]
La condition $\sigma_0x\ne\sigma_1x\quad\forall x\in R$ se transcrit par
\[
\bigm|\,u^{-1}x\ne x\qquad\forall x\in R
\]
donc ensemble avec les deux précédentes
cela signifie :
\[
\forall x\in R,\qquad\mathrm{Card}\{x,\sigma x,ux\}=3
\]
Noter que l'on a
\diagram{\begin{tikzcd}[ampersand replacement=\&,column sep=small]
1\arrow[r]\& \underbrace{G_1^{\unsure{0|\circ}}}_{\unsure{\simeq\mathbb Z|\mathbb Z}}\arrow[r]\& G_1\arrow[r]\&\deleted{\unsure{G}}\mathbb Z/2\mathbb Z\arrow[r]\&1
\end{tikzcd}}
Noter que $G_1\to G_1^{0}$ est formé d'él. d'ordre 2, \uline{chacun} définit splitting
% archive sheet number: 55
